\chapter{ACA-120 --- Verifiable Authority Cell}

\section{Scope}
This chapter defines the normative contract for a Verifiable Authority Cell
(VAC). A VAC is the smallest reusable ACA decision primitive that evaluates a
declared set of authority, jurisdiction, rule, evidence, procedure, time,
notice, conflict, and prohibition predicates and emits a canonical decision
state with sufficient proof and remedy metadata for independent inspection.

A conforming VAC does not determine whether an assertion is true merely because
it is well formed, signed, or hashed. It evaluates declared inputs under a
versioned profile and records the basis, limits, and consequences of that
evaluation.

\section{Conformance language}
The key words SHALL, SHALL NOT, SHOULD, SHOULD NOT, and MAY are to be
interpreted as normative requirements. A deployment claiming ACA-120
conformance SHALL identify the profile, schema version, rule version,
evaluator version, and conformance-test version used for each evaluation.

\begin{requirement}[ACA-120-001]
A conforming implementation SHALL satisfy every applicable SHALL and SHALL NOT
requirement in this chapter and SHALL produce evidence sufficient to verify that
claim.
\end{requirement}

\begin{requirement}[ACA-120-002]
A profile MAY add constraints or safe states, but it SHALL NOT weaken a baseline
safety requirement, erase provenance, reinterpret a canonical state, or permit
an action that the baseline contract requires to be held, denied, revoked, or
escalated.
\end{requirement}

\section{Authority Cell identity and boundary}
Each VAC evaluation SHALL have a unique cell identifier and evaluation
identifier. The cell identifier names the reusable decision primitive; the
evaluation identifier names one execution of that primitive.

\begin{requirement}[ACA-120-003]
The cell boundary SHALL declare its inputs, outputs, governing profile,
authoritative rule sources, required evidence classes, permitted callers,
permitted downstream actions, and remedy route.
\end{requirement}

\begin{requirement}[ACA-120-004]
Inputs not declared by the governing profile SHALL NOT silently influence the
evaluation. Any external human or automated influence SHALL be represented as a
versioned input, override, or review event.
\end{requirement}

\section{Canonical predicate model}
A baseline cell evaluates:
\[
I=\{A,J,R,E,P,T,N,C,V\}
\]
where $A$ is authority validity, $J$ jurisdiction, $R$ rule applicability, $E$
evidence sufficiency, $P$ required procedure, $T$ temporal validity, $N$ notice
and review, $C$ conflict state, and $V$ revocation or prohibition.

Each material predicate SHALL include:
\begin{enumerate}
\item a stable predicate identifier;
\item a declared type and value domain;
\item whether it is mandatory or advisory;
\item the rule or profile provision that requires it;
\item its evaluated value;
\item the evidence or authority references used;
\item the evaluator identity and evaluation time;
\item its confidence, uncertainty, or non-computability status where applicable;
\item an explanation code suitable for testing and review.
\end{enumerate}

\begin{requirement}[ACA-120-005]
Predicate values SHALL be drawn from a declared finite domain. The baseline
domain SHALL distinguish at least \state{TRUE}, \state{FALSE},
\state{UNKNOWN}, \state{CONFLICTED}, and \state{NOT-APPLICABLE}.
\end{requirement}

\begin{requirement}[ACA-120-006]
A mandatory predicate with value \state{UNKNOWN} SHALL prevent
\state{ALLOW}.
\end{requirement}

\begin{requirement}[ACA-120-007]
A mandatory predicate with value \state{CONFLICTED} SHALL produce
\state{CONFLICT} unless a profile defines a higher-priority safe state.
\end{requirement}

\begin{requirement}[ACA-120-008]
\state{NOT-APPLICABLE} SHALL be permitted only when the governing profile
contains an explicit applicability rule that can be independently evaluated.
It SHALL NOT be used as a substitute for missing evidence or unresolved scope.
\end{requirement}

\begin{requirement}[ACA-120-009]
Profiles SHALL identify which predicates are objective, evidentiary,
discretionary, or non-computable. Human judgment SHALL NOT be concealed as
deterministic computation.
\end{requirement}

\section{Authority and jurisdiction binding}
The authority object and jurisdiction predicate establish who may decide what,
for whom, where, during what period, and subject to which delegation,
prohibition, and review structure.

\begin{requirement}[ACA-120-010]
An authority reference SHALL identify the claimed source, issuer, delegate,
scope, affected subject or resource class, jurisdiction, effective period,
version, and revocation status.
\end{requirement}

\begin{requirement}[ACA-120-011]
A delegation SHALL NOT convey broader scope, longer duration, greater remedy
immunity, or higher precedence than the delegating authority possesses.
\end{requirement}

\begin{requirement}[ACA-120-012]
The subject, resource, transaction, or decision target evaluated by the cell
SHALL be explicitly bound to the authority and jurisdiction inputs. Name,
identifier, role, or record similarity alone SHALL NOT establish that binding.
\end{requirement}

\begin{requirement}[ACA-120-013]
Where authority or jurisdiction depends on a disputed fact, the dependency
SHALL be represented as a material predicate and SHALL NOT be assumed by the
evaluator.
\end{requirement}

\section{Rule and profile binding}
A VAC evaluates a declared rule set rather than an unversioned narrative.

\begin{requirement}[ACA-120-014]
Every evaluation SHALL bind to an immutable rule version and profile version.
A later rule change SHALL NOT retroactively alter the recorded result of an
earlier evaluation; it SHALL create a new evaluation or review event.
\end{requirement}

\begin{requirement}[ACA-120-015]
The profile SHALL declare predicate dependencies, mandatory predicates,
precedence rules, state-resolution rules, permitted overrides, and remedy
obligations in a form suitable for conformance testing.
\end{requirement}

\begin{requirement}[ACA-120-016]
Free-form prose MAY accompany a rule or predicate, but prose alone SHALL NOT be
treated as executable semantics when a conformance claim depends on that
semantics.
\end{requirement}

\section{Evidence sufficiency and provenance}
Evidence sufficiency is profile-relative and SHALL remain distinct from evidence
integrity and factual truth.

\begin{requirement}[ACA-120-017]
Each material evidence reference SHALL identify provenance, acquisition time,
source identity or source class, integrity metadata, access restrictions,
status, and the assertion it is offered to support or contradict.
\end{requirement}

\begin{requirement}[ACA-120-018]
A hash, signature, timestamp, or chain-of-custody record MAY support integrity
or attribution, but SHALL NOT by itself be treated as proof that the underlying
assertion is true, complete, lawful, or applicable.
\end{requirement}

\begin{requirement}[ACA-120-019]
The evaluator SHALL record material evidence that supports, contradicts, limits,
or remains unavailable for a predicate. Selective omission of known material
contradictory evidence is non-conforming.
\end{requirement}

\begin{requirement}[ACA-120-020]
Where lawful access restrictions prevent disclosure, the Decision Proof Record
SHALL preserve the existence, class, custodian, access rule, and effect of the
restricted evidence without exposing protected content.
\end{requirement}

\section{Procedure, notice, and review}
Procedure and notice are independently testable authorization conditions, not
post-decision annotations.

\begin{requirement}[ACA-120-021]
A required procedural step SHALL identify its responsible actor, triggering
condition, permitted sequence, completion evidence, deadline, and consequence
of non-completion.
\end{requirement}

\begin{requirement}[ACA-120-022]
Where notice or review is required before authorization or execution, absent,
untimely, defective, or unverified notice SHALL prevent \state{ALLOW} unless a
published profile defines a lawful emergency path and its additional safeguards.
\end{requirement}

\begin{requirement}[ACA-120-023]
An emergency path SHALL be explicit, time-bounded, independently reviewable, and
incapable of silently becoming the ordinary path.
\end{requirement}

\section{Temporal validity}
Time SHALL be evaluated as part of authority, rule, evidence, procedure, and
remedy validity.

\begin{requirement}[ACA-120-024]
The evaluator SHALL record the evaluation time, applicable effective intervals,
expiration conditions, and the clock or trusted time source used where time is
material.
\end{requirement}

\begin{requirement}[ACA-120-025]
Expired authority, expired evidence where freshness is mandatory, or an expired
procedural window SHALL prevent \state{ALLOW} unless the governing profile
contains an explicit and testable exception.
\end{requirement}

\section{Conflict, prohibition, and revocation precedence}
For profiles that permit Boolean authorization, the baseline allow expression
is:
\[
\state{ALLOW}=A\land J\land R\land E\land P\land T\land N\land \neg C\land \neg V.
\]
This expression is a minimum authorization condition, not a complete execution
instruction.

\begin{requirement}[ACA-120-026]
A valid revocation or prohibition SHALL take precedence over ordinary
authorization.
\end{requirement}

\begin{requirement}[ACA-120-027]
A conflict SHALL identify the conflicting assertions, authorities, rules, or
evidence objects and the precedence rule, review route, or missing fact required
to resolve it.
\end{requirement}

\begin{requirement}[ACA-120-028]
An implementation SHALL NOT resolve a conflict by discarding the lower-ranked
or inconvenient input from the proof record.
\end{requirement}

\section{Decision states and resolution}
A VAC SHALL emit one canonical state defined by ACA-130:
\[
S\in\{\state{ALLOW},\state{DENY},\state{HOLD},\state{CONFLICT},
\state{REVOKED},\state{ESCALATE}\}.
\]

\begin{requirement}[ACA-120-029]
\state{ALLOW} SHALL require every mandatory positive authorization predicate to
be satisfied and every mandatory conflict, revocation, or prohibition predicate
to be cleared.
\end{requirement}

\begin{requirement}[ACA-120-030]
\state{DENY} SHALL identify at least one evaluated rule condition that affirmatively
requires denial. Missing information alone SHALL ordinarily produce
\state{HOLD} or \state{ESCALATE}, not \state{DENY}.
\end{requirement}

\begin{requirement}[ACA-120-031]
\state{HOLD} SHALL identify the unresolved predicate, required evidence or
procedure, responsible route, and applicable deadline.
\end{requirement}

\begin{requirement}[ACA-120-032]
\state{ESCALATE} SHALL identify why the cell cannot safely or lawfully resolve
the evaluation and the qualified human or higher-order authority authorized to
review it.
\end{requirement}

\begin{requirement}[ACA-120-033]
Every emitted state SHALL include a stable reason code and a human-readable
explanation that is consistent with the recorded predicates and precedence
rules.
\end{requirement}

\section{Authorization and execution separation}
An authorization decision and the execution of an action are distinct events.

\begin{requirement}[ACA-120-034]
A VAC SHALL NOT represent an authorized decision as executed unless a separate
execution record identifies the action, executor, time, target, authorization
reference, result, and integrity mechanism.
\end{requirement}

\begin{requirement}[ACA-120-035]
Before execution, the executing component SHALL verify that the authorization
remains current, unrevoked, applicable to the target, and within any declared
execution window.
\end{requirement}

\begin{requirement}[ACA-120-036]
A failed, partial, delayed, duplicated, or reversed execution SHALL be recorded
without rewriting the original authorization decision.
\end{requirement}

\section{Decision Proof Record}
Every evaluation SHALL emit or reference a Decision Proof Record conforming to
ACA-140.

\begin{requirement}[ACA-120-037]
The proof record SHALL contain the cell and evaluation identifiers, profile and
rule versions, evaluator identity, material inputs, predicate results, emitted
state, reason code, conflicts and unknowns, authority and evidence references,
remedy route, and integrity metadata.
\end{requirement}

\begin{requirement}[ACA-120-038]
The proof record SHALL preserve sufficient normalized input and configuration
metadata to permit deterministic replay when the evaluator has lawful access to
the same versions and evidence objects.
\end{requirement}

\begin{requirement}[ACA-120-039]
The record SHALL distinguish an integrity proof, an attribution proof, an
authority evaluation, and a factual truth claim. None SHALL be silently
substituted for another.
\end{requirement}

\section{Replay and determinism}

\begin{requirement}[ACA-120-040]
A conforming implementation SHALL declare its canonicalization rules for object
ordering, numeric representation, text normalization, time representation,
identifier comparison, omitted values, and cryptographic digest construction.
\end{requirement}

\begin{requirement}[ACA-120-041]
Given the same normalized inputs, rule version, profile version, evaluator
version, and permitted deterministic dependencies, replay SHALL produce the
same predicate values, state, and reason code.
\end{requirement}

\begin{requirement}[ACA-120-042]
Non-deterministic services, probabilistic models, or mutable external data MAY
inform a predicate only when their identity, version, input, output, uncertainty,
and review rule are recorded. Their output SHALL NOT be mislabeled as
deterministic replay evidence.
\end{requirement}

\section{Remedy attachment}
A terminal or adverse state does not complete the ACA lifecycle unless the
applicable remedy is defined.

\begin{requirement}[ACA-120-043]
Every \state{DENY}, \state{HOLD}, \state{CONFLICT}, \state{REVOKED}, or
\state{ESCALATE} state SHALL identify an available remedy or shall explicitly
record that no remedy is provided, by what authority, and subject to what review.
\end{requirement}

\begin{requirement}[ACA-120-044]
A remedy reference SHALL identify the responsible authority, permitted filer or
initiator, required inputs, deadline, review sequence, correction mechanism,
propagation obligations, verification condition, and closure rule.
\end{requirement}

\begin{requirement}[ACA-120-045]
A later correction SHALL preserve the original evaluation and create a linked
correction, propagation, and verification history conforming to ACA-150.
\end{requirement}

\section{Security and abuse resistance}

\begin{requirement}[ACA-120-046]
Implementations SHALL enforce least privilege for evaluation, review, override,
execution, evidence access, and remedy closure. Possession of write access to
one function SHALL NOT imply authority to perform the others.
\end{requirement}

\begin{requirement}[ACA-120-047]
Overrides SHALL identify the authorized actor, legal or policy basis, scope,
reason, time, affected predicates, resulting state, expiration, and required
independent review. Undocumented override behavior is non-conforming.
\end{requirement}

\begin{requirement}[ACA-120-048]
The implementation SHALL detect or prevent replay substitution, stale authority,
identifier collision, unauthorized profile changes, evidence-reference
replacement, and suppression of material conflicts to the extent required by
the deployment threat model.
\end{requirement}

\begin{requirement}[ACA-120-049]
Sensitive evidence and personal data SHALL be minimized and access-controlled,
but minimization SHALL NOT erase the proof that a material input existed or
alter its recorded effect on the decision.
\end{requirement}

\section{Conformance evidence and traceability}

\begin{requirement}[ACA-120-050]
Each normative requirement claimed by an implementation SHALL map to at least
one specification reference, implementation control, positive test, negative or
adversarial test where applicable, and retained test result.
\end{requirement}

\begin{requirement}[ACA-120-051]
A schema-valid object SHALL NOT be labeled substantively conforming unless its
cross-object bindings, predicate semantics, state resolution, proof record,
replay behavior, and remedy obligations also pass the applicable conformance
suite.
\end{requirement}

\begin{requirement}[ACA-120-052]
A conformance report SHALL identify every requirement tested, passed, failed,
not applicable, or not tested. Partial conformance SHALL be labeled as partial
and SHALL NOT be presented as full ACA-120 conformance.
\end{requirement}

\section{Versioning and extension rules}

\begin{requirement}[ACA-120-053]
Every extension predicate, state restriction, evidence class, or precedence rule
SHALL use a collision-resistant identifier and SHALL declare its relationship to
the baseline contract.
\end{requirement}

\begin{requirement}[ACA-120-054]
A breaking semantic change SHALL require a new contract or profile version and
SHALL NOT be introduced solely by changing prose, examples, evaluator code, or
an external service.
\end{requirement}

\begin{requirement}[ACA-120-055]
Deprecated fields or rules SHALL remain interpretable for the declared retention
period, and migrations SHALL preserve the original version, meaning, proof
record, and remedy history.
\end{requirement}

\section{Minimum acceptance criteria}
An ACA-120 candidate is eligible for approval only when:
\begin{enumerate}
\item the normative text is versioned and reviewable;
\item the Authority Cell schema expresses the required structural contract;
\item positive, negative, conflict, revocation, expiration, unknown, override,
and replay examples exist;
\item a requirement-to-schema-to-test traceability matrix exists;
\item the conformance suite verifies state resolution and fail-safe behavior;
\item authorization and execution are represented separately;
\item proof-record and remedy obligations are tested; and
\item all known deviations are documented as explicit, time-bounded exceptions.
\end{enumerate}

This chapter defines the normative Authority Cell contract. ACA-130 defines the
state-machine transitions, ACA-140 defines the portable Decision Proof Record,
and ACA-150 defines the remedy lifecycle that follows an unresolved, adverse,
revoked, or corrected decision.